\chapter{The RISC-V Processor: Architecture, Implementation, and Verification}\label{riscvProcessor}
\index{RISC-V: open standard instruction set architecture}\index{RISC-V!architecture, implementation, and verification}
\localtoc

%*******************************************************************
\section{Introduction to the RISC-V Architecture}\label{riscvIntro}
%*******************************************************************
\index{RISC-V!design philosophy}\index{open standard instruction set}

The RISC-V (pronounced ``risk-five'') instruction set architecture represents a fundamental paradigm shift in computer architecture. Unlike proprietary instruction set architectures (such as x86 and ARM) whose definitions are tightly controlled by commercial entities, RISC-V is an open, royalty-free standard governed by RISC-V International. Conceived in 2010 at the University of California, Berkeley \cite{asanovic2014case,waterman2014riscv,waterman2016design,patterson2017riscvreader}, RISC-V was developed to provide a clean, modern, extensible, and mathematically structured instruction set architecture suitable for academic instruction, research, and high-performance commercial silicon. The formal definition of the architecture is maintained in the RISC-V Instruction Set Manual, comprising Volume I for the Unprivileged ISA \cite{riscv-spec}, Volume II for the Privileged Architecture \cite{riscv-priv-spec}, and modular ratified extension standards including the \texttt{Zmmul} multiplication extension \cite{riscv-zmmul-spec}.

The architectural philosophy of RISC-V is guided by four foundational principles:
\begin{enumerate}
    \item \textbf{Simplicity and Regularity}: The instruction set is structured around a small, base integer instruction set (\texttt{RV32I} for 32-bit addresses, \texttt{RV64I} for 64-bit addresses). Instructions maintain a uniform 32-bit length, fixed opcode locations, and symmetric register source and destination fields.
    \item \textbf{Strict Modularity}: Rather than obligating every implementation to support an expansive, monolithic feature set, RISC-V organizes capabilities into modular standard extensions. A minimalist embedded core may implement only the base integer instructions (\texttt{RV32I}), while an accelerator or application processor can incrementally add standard extensions for multiplication and division (\texttt{M}), atomic memory operations (\texttt{A}), single-precision floating-point (\texttt{F}), double-precision floating-point (\texttt{D}), and vector processing (\texttt{V}).
    \item \textbf{Clear Separation of Privileged and Unprivileged Architecture}: Instruction execution, computational semantics, and user-level register manipulation are strictly decoupled from machine-mode control, exception management, physical memory protection, and virtual memory translation.
    \item \textbf{Extensibility Without Fragmentation}: The base instruction set reserves dedicated opcode spaces and major instruction formats specifically for custom, application-specific instructions without conflicting with standardized extensions.
\end{enumerate}

\subsection{Architectural State and Programmer's Model}
\index{RISC-V!programmer's model}\index{register file!RISC-V architectural state}

The unprivileged programmer's model for the 32-bit base integer architecture (\texttt{RV32I}) consists of:
\begin{itemize}
    \item \textbf{Thirty-two General-Purpose Registers ($x0$--$x31$)}: Each register has a width of $\text{XLEN} = 32$ bits. Register $x0$ is hardwired to the constant zero ($x0 \equiv 0$) on both read and write operations; writes to $x0$ have no effect, and reading $x0$ always yields zero. Registers $x1$ through $x31$ are general-purpose registers that hold integer values, pointers, and temporary evaluation results.
    \item \textbf{Program Counter (\texttt{pc})}: A 32-bit register holding the byte address of the instruction currently being fetched and executed.
    \item \textbf{Linear Byte-Addressable Memory}: The address space comprises $2^{32}$ bytes, organized in a little-endian byte ordering wherein the least significant byte of a multi-byte word resides at the lowest memory address.
\end{itemize}

Table \ref{tab:riscv_abi_registers} summarizes the standard RISC-V register conventions (Application Binary Interface, or ABI).

\begin{table}[htbp]
\centering
\small
\begin{tabular}{|l|l|l|l|}
\hline
\textbf{Register} & \textbf{ABI Name} & \textbf{Description} & \textbf{Saver} \\
\hline
\texttt{x0}       & \texttt{zero}     & Hardwired zero & --- \\
\texttt{x1}       & \texttt{ra}       & Return address & Caller \\
\texttt{x2}       & \texttt{sp}       & Stack pointer & Callee \\
\texttt{x3}       & \texttt{gp}       & Global pointer & --- \\
\texttt{x4}       & \texttt{tp}       & Thread pointer & --- \\
\texttt{x5}--\texttt{x7}   & \texttt{t0}--\texttt{t2}   & Temporaries & Caller \\
\texttt{x8}       & \texttt{s0}/\texttt{fp} & Saved register / Frame pointer & Callee \\
\texttt{x9}       & \texttt{s1}       & Saved register & Callee \\
\texttt{x10}--\texttt{x11} & \texttt{a0}--\texttt{a1}   & Function arguments / Return values & Caller \\
\texttt{x12}--\texttt{x17} & \texttt{a2}--\texttt{a7}   & Function arguments & Caller \\
\texttt{x18}--\texttt{x27} & \texttt{s2}--\texttt{s11}  & Saved registers & Callee \\
\texttt{x28}--\texttt{x31} & \texttt{t3}--\texttt{t6}   & Temporaries & Caller \\
\hline
\end{tabular}
\caption{RISC-V unprivileged general-purpose registers and standard ABI roles.}\label{tab:riscv_abi_registers}
\end{table}


\section{The Instruction Set Architecture: RV32I and Zmmul}\label{riscvISA}
\index{RV32I: RISC-V base 32-bit integer ISA}\index{Zmmul: standard multiplication extension}

\subsection{Instruction Formats}
\index{RISC-V!instruction formats}\index{instruction encoding!RISC-V formats}

All RISC-V base instructions are exactly 32 bits in width and are aligned on 4-byte memory boundaries. To maximize decoding speed and eliminate unnecessary multiplexers in high-performance hardware, RISC-V employs six primary instruction formats (Figure \ref{fig:riscv_formats}):

\begin{figure}[htbp]
\centering
\small
\begin{tabular}{|c|c|c|c|c|c|c|}
\hline
\textbf{Format} & \textbf{Bits 31:25} & \textbf{Bits 24:20} & \textbf{Bits 19:15} & \textbf{Bits 14:12} & \textbf{Bits 11:7} & \textbf{Bits 6:0} \\
\hline
\textbf{R-type} & \texttt{funct7} (7) & \texttt{rs2} (5) & \texttt{rs1} (5) & \texttt{funct3} (3) & \texttt{rd} (5) & \texttt{opcode} (7) \\
\textbf{I-type} & \multicolumn{2}{c|}{\texttt{imm[11:0]} (12)} & \texttt{rs1} (5) & \texttt{funct3} (3) & \texttt{rd} (5) & \texttt{opcode} (7) \\
\textbf{S-type} & \texttt{imm[11:5]} (7) & \texttt{rs2} (5) & \texttt{rs1} (5) & \texttt{funct3} (3) & \texttt{imm[4:0]} (5) & \texttt{opcode} (7) \\
\textbf{B-type} & \texttt{imm[12|10:5]} (7) & \texttt{rs2} (5) & \texttt{rs1} (5) & \texttt{funct3} (3) & \texttt{imm[4:1|11]} (5) & \texttt{opcode} (7) \\
\textbf{U-type} & \multicolumn{4}{c|}{\texttt{imm[31:12]} (20)} & \texttt{rd} (5) & \texttt{opcode} (7) \\
\textbf{J-type} & \multicolumn{4}{c|}{\texttt{imm[20|10:1|11|19:12]} (20)} & \texttt{rd} (5) & \texttt{opcode} (7) \\
\hline
\end{tabular}
\caption{The six fundamental RISC-V instruction formats.}\label{fig:riscv_formats}
\end{figure}

Notice the structural regularity across all six formats:
\begin{itemize}
    \item The major 7-bit opcode field ($\text{inst}[6:0]$) is always positioned in bits 6 through 0.
    \item Destination register address ($\text{rd}$) is always located in bits 11 through 7.
    \item First source register address ($\text{rs1}$) is always located in bits 19 through 15.
    \item Second source register address ($\text{rs2}$) is always located in bits 24 through 20.
    \item The 3-bit sub-operation specifier ($\text{funct3}$) is always located in bits 14 through 12.
    \item The sign bit of all immediate values is unconditionally mapped to bit 31 ($\text{inst}[31]$), allowing sign-extension hardware to operate concurrently with instruction decoding.
\end{itemize}

\subsection{Instruction Classes}
\index{RISC-V!instruction classes}

The \texttt{RV32I} base instruction set comprises 40 machine instructions categorized into four functional groups:

\subsubsection{Register-Register and Register-Immediate Computational Operations}
\index{ALU!RISC-V operations}\index{computational instructions}
\begin{itemize}
    \item \textbf{Arithmetic}: Addition (\texttt{ADD}, \texttt{ADDI}) and Subtraction (\texttt{SUB}). Unsigned arithmetic is performed using identical two's complement adder circuits.
    \item \textbf{Logical}: Bitwise Boolean operations AND (\texttt{AND}, \texttt{ANDI}), OR (\texttt{OR}, \texttt{ORI}), and exclusive-OR (\texttt{XOR}, \texttt{XORI}).
    \item \textbf{Shifts}: Logical Shift Left (\texttt{SLL}, \texttt{SLLI}), Logical Shift Right (\texttt{SRL}, \texttt{SRLI}), and Arithmetic Shift Right (\texttt{SRA}, \texttt{SRAI}). In 32-bit mode, the shift amount is masked to the lower 5 bits ($\text{shamt}[4:0]$).
    \item \textbf{Comparisons}: Set Less Than signed (\texttt{SLT}, \texttt{SLTI}) and Set Less Than unsigned (\texttt{SLTU}, \texttt{SLTIU}), writing 1 to \texttt{rd} if $\text{rs1} < \text{rs2}$ (or $\text{rs1} < \text{imm}$) and 0 otherwise.
    \item \textbf{Upper Immediates}: Load Upper Immediate (\texttt{LUI}), which places a 20-bit immediate in bits 31:12 of \texttt{rd}, and Add Upper Immediate to PC (\texttt{AUIPC}), which computes $\text{pc} + (\text{imm} \ll 12)$ for position-independent addressing.
\end{itemize}

\subsubsection{Control Transfer Instructions}
\index{branch instructions!RISC-V}\index{jump instructions!JAL and JALR}
\begin{itemize}
    \item \textbf{Conditional Branches (B-type)}: Branch if Equal (\texttt{BEQ}), Branch if Not Equal (\texttt{BNE}), Branch if Less Than signed (\texttt{BLT}), Branch if Greater Than or Equal signed (\texttt{BGE}), Branch if Less Than unsigned (\texttt{BLTU}), and Branch if Greater Than or Equal unsigned (\texttt{BGEU}). The target address is computed as $\text{pc} + \text{imm}$, where the 12-bit signed immediate represents a multiple of 2 bytes with bit 0 implicitly 0.
    \item \textbf{Unconditional Jumps (J-type and I-type)}: Jump and Link (\texttt{JAL}) computes target address $\text{pc} + \text{imm}$ and writes $\text{pc} + 4$ to destination register \texttt{rd}. Jump and Link Register (\texttt{JALR}) computes target address $(\text{rs1} + \text{imm})$ and writes $\text{pc} + 4$ to \texttt{rd}. Per the RISC-V specification, the least significant bit of the JALR target is explicitly cleared to zero:
    \begin{equation}
        \text{target}_{\text{JALR}} = (\text{rs1} + \text{imm}) \ \& \ \sim 1
    \end{equation}
\end{itemize}

\subsubsection{Memory Load and Store Instructions}
\index{memory!RISC-V load and store}\index{byte addressing}
\begin{itemize}
    \item \textbf{Loads (I-type)}: Load Byte signed (\texttt{LB}), Load Halfword signed (\texttt{LH}), Load Word (\texttt{LW}), Load Byte unsigned (\texttt{LBU}), and Load Halfword unsigned (\texttt{LHU}). The memory address is computed as $\text{rs1} + \text{imm}$. Unsigned variants zero-extend the sub-word data, while signed variants sign-extend to 32 bits.
    \item \textbf{Stores (S-type)}: Store Byte (\texttt{SB}), Store Halfword (\texttt{SH}), and Store Word (\texttt{SW}). Data from source register \texttt{rs2} is written to effective address $\text{rs1} + \text{imm}$.
\end{itemize}

\subsubsection{The Zmmul Standard Extension}
\index{Zmmul: multiplication extension}\index{hardware multiplier!RISC-V}
The standard \texttt{Zmmul} extension defines the integer multiplication subset of the standard M extension without requiring hardware division circuitry:
\begin{itemize}
    \item \texttt{MUL}: Computes the low 32 bits of $\text{rs1} \times \text{rs2}$.
    \item \texttt{MULH}: Computes the high 32 bits of signed $\times$ signed product $\text{rs1} \times \text{rs2}$.
    \item \texttt{MULHSU}: Computes the high 32 bits of signed $\times$ unsigned product $\text{rs1} \times \text{rs2}$.
    \item \texttt{MULHU}: Computes the high 32 bits of unsigned $\times$ unsigned product $\text{rs1} \times \text{rs2}$.
\end{itemize}


%***********************************************************************************
\subsection{Harvard 5-OS Hardware Architecture and Microarchitecture}\label{riscvArch}
%***********************************************************************************
\index{Harvard architecture!RISC-V 5-OS}\index{processor architecture!5-OS loops}

As described in the structural classification of computing systems Chapter \ref{lect6a}, our RISC-V processor is implemented as a \textbf{5th-Order System (Harvard 5-OS)} comprising five coupled structural feedback loops:
\begin{enumerate}
    \item \textbf{1st loop (1-OS) --- Memory State Loop}: Closed by the architectural register file (\texttt{RiscvRegFile}). Source operands read from \texttt{rs1} and \texttt{rs2} are processed and rewritten synchronously to destination register \texttt{rd} on the rising clock edge when \texttt{regWrite} is asserted, preserving state across instruction steps. Register $x0$ is hardwired to zero on both read and write.
    \item \textbf{2nd loop (2-OS) --- Autonomous Subsystem Loops}: Operates two internal autonomous automata:
    \begin{itemize}
        \item The \textbf{Program Counter Automaton} (\texttt{ProgramCounter}), updating execution address $\text{pc} \leftarrow \text{pc} + 4$.
        \item The \textbf{Registers with Arithmetic Logic Unit} (\texttt{RALU}), closing the internal computational feedback loop between the register file and the \texttt{RiscvALU}.
    \end{itemize}
    \item \textbf{3rd loop (3-OS) --- Executive-Control Loop}: Represents the interactive coupling between the RALU execution subsystem and the Program Counter automaton:
    \begin{itemize}
        \item The \texttt{RiscvALU} computes relational flags (\texttt{zeroFlag}, \texttt{lessThanFlag}, \texttt{lessThanUFlag}) that drive branch condition multiplexers, dynamically updating the PC automaton when branch conditions evaluate true.
        \item The PC automaton routes the subroutine link address ($\text{pc} + 4$) to the RALU write-back path to be preserved in link register \texttt{ra} ($x1$) for call sequences (\texttt{JAL}, \texttt{JALR}).
    \end{itemize}
    \item \textbf{4th loop (4-OS) --- Program Memory Loop}: The PC automaton presents execution address \texttt{io.imem.addr} to the Program Memory. Instruction word \texttt{io.imem.inst} is delivered combinationally to the instruction decoder (\texttt{RiscvDecoder}) to configure datapath multiplexers.
    \item \textbf{5th loop (5-OS) --- Data Memory Loop}: The RALU drives computed effective address $\text{rs1} + \text{imm}$ and store data \texttt{rs2Data} to byte-addressable memory (\texttt{RiscvDataMemory}), and receives load data into the register write-back multiplexer.
\end{enumerate}

%---------------------------------------------
\subsubsection{Single-Cycle Fetch-Execute Cycle}
%----------------------------------------------
\index{Fetch-Execute cycle!single-cycle RISC-V}

A single-cycle fetch-execute processor executes each instruction within a single clock cycle across five coordinated logical functions:
\begin{enumerate}
    \item \textbf{Instruction Fetch (IF)}: The \texttt{ProgramCounter} presents address \texttt{pc.io.pc} onto \texttt{io.imem.addr}. Unbuffered memory combinationally returns the 32-bit instruction word \texttt{inst}.
    \item \textbf{Instruction Decode (ID)}: \texttt{RiscvDecoder} unpacks register addresses (\texttt{rd}, \texttt{rs1}, \texttt{rs2}), while \texttt{RiscvImmGen} produces sign-extended 32-bit immediate \texttt{d.imm}. Direct-wire bit extraction maps $\{\text{inst}[25], \text{inst}[30], \text{inst}[14:12]\}$ directly into \texttt{c.aluOp}.
    \item \textbf{Operand Fetch \& Execution (EX)}: \texttt{RALU} reads operands \texttt{rs1Data} and \texttt{rs2Data}. Multiplexers select between \texttt{rs2Data} and \texttt{d.imm} based on \texttt{c.aluSrc}. \texttt{RiscvALU} computes the 32-bit result and evaluation flags.
    \item \textbf{Memory Access \& Write-Back (MEM/WB)}: Memory access controls (\texttt{memRead}, \texttt{memWrite}, \texttt{funct3}) are driven to \texttt{io.dmem}. The write-back multiplexer selects between \texttt{aluResult}, memory load data, $\text{pc} + 4$ (for jumps), and upper immediate constants, writing to destination register \texttt{rd} on the rising clock edge.
    \item \textbf{Next PC Generation (PC)}: When a branch condition is satisfied or a jump is decoded, target address $\text{pc} + \text{imm}$ or $\text{rs1} + \text{imm}$ is loaded into the PC automaton via \texttt{Mode.Jalr} (enforcing LSB $= 0$). Otherwise, the PC automaton increments by $+4$.
\end{enumerate}

\subsubsection{Chisel Hardware Implementation}
\index{Chisel!RiscvFetchExecute implementation}

Listing \ref{lst:chisel_riscv_core_ch18} presents the complete Chisel implementation of the single-cycle processor core.

\includechiselcode[caption={RISC-V RV32I\_Zmmul 5-OS Fetch-Execute Processor Core (Chisel: RiscvFetchExecute.scala)}, label={lst:chisel_riscv_core_ch18}, firstline=10, firstnumber=10, lastline=125]{\chiselSourceDir/riscv/RiscvFetchExecute.scala}


%***********************************************************************
\section{Architectural Conformance Verification}\label{riscvConformance}
%************************************************************************
\index{conformance testing!RISC-V architectural tests}\index{verification!Spike co-simulation}

A challenge in processor engineering is guaranteeing that a hardware implementation strictly conforms to the formal instruction set specification. While ad-hoc assembly tests can verify basic functionality, they rarely test corner cases such as arithmetic overflow boundary conditions, sign extension edge cases, misaligned immediate offsets, and bit-level side effects.

To establish correctness, all processor implementations in this chapter are verified against the official RISC-V Architectural Conformance Suite (\texttt{riscv-arch-test}) developed by RISC-V International \cite{riscv-arch-test}.

\subsection{Verification Methodology: Pure Chisel/Scala Simulation}
\index{EphemeralSimulator!architectural testbench}

While we can generate Verilog and compile a C++ simulator using external tools (such as Verilator), we choose to use Chisel's built-in simulator (\texttt{EphemeralSimulator}) 
\footnote{For the prior approach, see the KryptoNyte processor environment. \url{https://github.com/Ypologist/KryptoNyte}}.
Our Ephemeral simulator approach achieves cycle-accurate software simulation entirely within the \texttt{sbt} test framework without requiring external synthesis or toolchain dependencies:
\begin{itemize}
    \item \textbf{4MB Emulated Memory System}: A simulation memory buffer is initialized at base address \texttt{0x80000000} (the canonical RISC-V boot vector).
    \item \textbf{Mailbox Completion Protocol}: The testbench continuously executes instructions and monitors data memory writes. When the processor writes to the memory-mapped mailbox register \texttt{tohost} (\texttt{0x80004000}), the test completes. A write value of \texttt{0x1} signifies success, whereas any other non-zero value indicates a failure code.
    \item \textbf{Golden Reference Signature Comparison}: Upon test completion, memory contents between linker symbols \texttt{begin\_signature} and \texttt{end\_signature} are extracted word-by-word into a test signature. This signature is compared against the golden reference signature produced by \textit{Spike}, the official RISC-V reference golden model \cite{riscv-spike}\footnote{\url{https://github.com/riscv-software-src/riscv-isa-sim}}.
\end{itemize}

Listing \ref{lst:chisel_conformance_spec} illustrates a partial listing of the Chisel/Scala conformance test runner (showing lines 94--111).

\includechiselcode[caption={Chisel/Scala Conformance Test Suite (RiscvConformanceSpec.scala, partial listing)}, label={lst:chisel_conformance_spec}, firstline=94, firstnumber=94, lastline=111]{\chiselTestBenchSourceDir/riscv/RiscvConformanceSpec.scala}

\subsection{Conformance Results}
\index{conformance testing!results summary}\index{conformance testing!RV32I verification}

All processor core configurations developed throughout Chapter 19 pass 100\% of the official architectural conformance tests with bit-exact signature matches against the official Spike reference golden model \cite{riscv-spike} across all test suites (spanning 38 suites for pure \texttt{RV32I} cores and 42 suites when the \texttt{Zmmul} extension is enabled). Table \ref{tab:conformance_results} details the complete suite of verified architectural test cases by instruction category. Detailed cycle-by-cycle execution metrics and comparative analysis for each core microarchitecture are presented in their respective results sections throughout this chapter, culminating in a grand verification summary in Section \ref{riscvSummary}.

\begin{table}[htbp]
\centering
\small
\begin{tabularx}{\textwidth}{|l|X|}
\hline
\textbf{Category} & \textbf{Architectural Test Suites (All Pass)} \\
\hline
\textbf{Arithmetic (R-type)} & \texttt{add-01}, \texttt{sub-01}, \texttt{sll-01}, \texttt{slt-01}, \texttt{sltu-01}, \texttt{xor-01}, \texttt{srl-01}, \texttt{sra-01}, \texttt{or-01}, \texttt{and-01} \\
\hline
\textbf{Immediate (I-type)}  & \texttt{addi-01}, \texttt{slli-01}, \texttt{slti-01}, \texttt{sltiu-01}, \texttt{xori-01}, \texttt{srli-01}, \texttt{srai-01}, \texttt{ori-01}, \texttt{andi-01} \\
\hline
\textbf{Upper Immediates}     & \texttt{lui-01}, \texttt{auipc-01} \\
\hline
\textbf{Jumps}               & \texttt{jal-01}, \texttt{jalr-01} \\
\hline
\textbf{Branches (B-type)}   & \texttt{beq-01}, \texttt{bne-01}, \texttt{blt-01}, \texttt{bge-01}, \texttt{bltu-01}, \texttt{bgeu-01} \\
\hline
\textbf{Loads (I-type)}      & \texttt{lb-align-01}, \texttt{lbu-align-01}, \texttt{lh-align-01}, \texttt{lhu-align-01}, \texttt{lw-align-01} \\
\hline
\textbf{Stores (S-type)}     & \texttt{sb-align-01}, \texttt{sh-align-01}, \texttt{sw-align-01} \\
\hline
\textbf{Synchronization}     & \texttt{fence-01} \\
\hline
\textbf{Zmmul Multiplication}& \texttt{mul-01}, \texttt{mulh-01}, \texttt{mulhsu-01}, \texttt{mulhu-01} \\
\hline
\multicolumn{2}{|l|}{\textbf{Total / Summary}: 42 official test suites (38 RV32I + 4 Zmmul) --- 100\% signature match across all cores} \\
\hline
\end{tabularx}
\caption{Official RISC-V architectural conformance verification test suites.}\label{tab:conformance_results}
\end{table}


\section{Embedded Benchmark Workloads and Performance Evaluation Methodology}\label{riscvBenchmarks}
\index{benchmark!workloads}\index{benchmark!taxonomy}\index{CoreMark}\index{Dhrystone}\index{Embench-IoT}

While architectural conformance suites verify functional correctness, they are not representative of real-world application performance. Microprocessor performance evaluation requires standardized benchmark workloads that exercise computational, memory, and control-flow subsystems in realistic combinations. In this chapter, we evaluate all processor configurations against four distinct benchmark workloads:

\begin{enumerate}
    \item \textbf{Digital Signal Processing (DSP) Kernels}:
    Signal processing workloads demand sustained multiply-accumulate throughput. We implement two canonical DSP kernels compiled as a bare-metal application:
    \begin{itemize}
        \item A \textbf{16-point 4-tap Finite Impulse Response (FIR) filter}:
        \begin{equation}
        y[n] = \sum_{k=0}^{3} h[k] \cdot x[n-k].
        \end{equation}
        \item A \textbf{16-element vector dot product}:
        \begin{equation}
        P = \sum_{i=0}^{15} a[i] \cdot b[i].
        \end{equation}
    \end{itemize}
    In pure \texttt{RV32I}, multiplications are executed using compiler runtime shift-and-add software subroutines (\texttt{\_\_mulsi3}), which require approximately 35 instructions per multiply. When the \texttt{Zmmul} extension is present, each multiplication is executed in a single hardware \texttt{MUL} instruction.
    
    \item \textbf{Dhrystone 2.1}:
    Originally developed by Reinhold P. Weicker in 1984 \cite{weicker1984dhrystone}, Dhrystone is a synthetic integer benchmark that models general-purpose system programming workloads. It contains representative mixes of integer arithmetic, string operations (\texttt{strcpy}, \texttt{strcmp}), pointer chasing, structure copying, and nested procedure calls. Dhrystone reports performance in Dhrystones per second, normalized against the DEC VAX-11/780 (rated at 1,757 Dhrystones/second = 1.0 DMIPS), yielding:
    \begin{equation}
    \text{DMIPS} = \frac{\text{Dhrystones/second}}{1757}, \qquad \text{DMIPS/MHz} = \frac{\text{DMIPS}}{f_{\text{clk}} \ (\text{MHz})}.
    \end{equation}
    
    \item \textbf{EEMBC CoreMark 1.0}:
    Developed by the Embedded Microprocessor Benchmark Consortium (EEMBC) \cite{galon2012coremark}, CoreMark is the de-facto industry standard for evaluating embedded microcontroller cores. Unlike synthetic benchmarks, CoreMark exercises algorithms commonly found in embedded systems: 2D matrix manipulation, linked-list search and reversal, state machine parsing, and a 16-bit Cyclic Redundancy Check (CRC-16). CoreMark incorporates a self-checking CRC that validates computational correctness, ensuring optimizing compilers cannot eliminate algorithmic work through dead-code removal. Performance is reported in:
    \begin{equation}
    \text{CoreMark/MHz} = \frac{\text{CoreMarks}}{f_{\text{clk}} \ (\text{MHz})} = \frac{10^6}{\text{Cycles per Iteration}}.
    \end{equation}
    
    \item \textbf{Embench-IoT}:
    Developed by the Free and Open Source Silicon (FOSSi) Foundation in collaboration with David Patterson \cite{embench2020}, Embench is a modern, open-source benchmark suite consisting of 19 real-world embedded programs (e.g., AES encryption, Huffman decoding, Minver, QR-code, Wic) designed specifically to provide a reproducible, non-proprietary evaluation standard for microcontroller-class RISC-V cores.
\end{enumerate}

Table \ref{tab:benchmark_taxonomy} summarizes the taxonomy of embedded benchmark standards used in this text.

\begin{table}[htbp]
\centering
\small
\resizebox{\linewidth}{!}{%
\begin{tabular}{|l|c|c|c|l|}
\hline
\textbf{Benchmark Suite} & \textbf{Primary Target} & \textbf{Runtime Environment} & \textbf{Footprint} & \textbf{Primary Metric} \\
\hline
\textbf{DSP Kernels}     & Embedded Signal Processing & Freestanding Bare-Metal & $<2$\,KB RAM & Cycles, Execution Time \\
\textbf{Dhrystone 2.1}   & General Integer Systems  & Freestanding Bare-Metal & $<4$\,KB RAM & DMIPS, DMIPS/MHz \\
\textbf{CoreMark 1.0}    & Embedded Microcontrollers & Freestanding Bare-Metal & $<8$\,KB RAM & CoreMark, CoreMark/MHz \\
\textbf{Embench-IoT}     & Embedded IoT Cores      & Freestanding Bare-Metal & $<64$\,KB RAM & Geometric Mean Relative Speed \\
\hline
\end{tabular}%
}
\caption{Taxonomy of embedded computer architecture benchmarks.}\label{tab:benchmark_taxonomy}
\end{table}

Throughout this chapter, rather than presenting benchmark data in isolation, results are reported in a core-by-core manner as each processor microarchitecture is developed: from the single-cycle baseline, through the unforwarded pipeline, balanced adder stages, data bypassing, pipelined multiplication, branch prediction, and compiler optimization flags.